% !TeX root = ../main.tex

\chapter{讨论与研究局限}
\section{四条路径的经济含义}
路径一保留持续成长与相对低估值的筛选骨架，以允许一次历史非正增长、相对排名和多维排序缓解候选稀缺。其全期收益高于E，但多项规则同步变化，不能把全部差异归于“19/20”或某一个因子。

路径二的价值在于恢复基础版可解释的全截面比较，同时展示排名分母如何进入组合结果；它在共同复现区间较优，在完整区间收益低于P1、回撤较小。规则更接近基础排名含义，不意味着历史业绩一定更高。

路径三、四保留各自选股名单，通过历史波动调整暴露，在当前数据中改善夏普和峰谷损失，但牺牲收益。P3全期指标较优，不应据此再叠加未验证的新规则、重新命名为唯一最优策略；四路径应共同保留，便于后续冻结检验。

\section{信息时点与参数选择}
最重要的限制是季度末财务映射。四路径的排名和风控修正均未解决财报实际披露时间和历史数值修订问题。此前冻结的公告敏感性结果显示，2003—2024年报告期映射年化25.40\%，完整公告约束为0.44\%；覆盖匹配仍为24.85\%。附录B保留其区间与定义，不将差值认定为前视偏差净因果效应，也不把该结果伪装成四路径已完成公告约束检验。

19/20及原有得分参数曾参考后续历史表现，P1与P2的保留选择也发生在研究者已见数据之后。追加区间是历史延长，不是严格事先冻结、从未接触的测试集。本文不报告未经校正的寻优显著性，也不把历史赢家排序表述为预测能力证据。

\section{收益、交易与统计边界}
收益采用含红利个股月收益与上证综指日收益复利，权重结构和分红口径并非完全一致。相对上证领先只是所设基准下的比较，不等于完成风险调整或证明可实现alpha。最大回撤基于月末，可能小于月内真实损失。

缺失收益再平均、月度重新等权、季度固定费用以及风控费用随仓位缩放，均是实现近似。未完整模拟涨跌停、停复牌、退市处置、冲击成本和融资约束。核心优先保留也不是当前组合逻辑。这些边界应在使用结果时与收益数字同时呈现。

\section{后续核验的优先顺序}
后续应先补证据，再增加策略复杂度。首先在CSMAR内部核查缺失公司的上市前比较期资料、合并报表类型和历史版本，分别记录“无报告期”“无同比基数”“基数非正”，避免将三者统称亏损。其次按原研报证券范围核对公告表和共同调仓日，只有事件时钟核对后，个股日行情验证才能针对真实差异。

若取得足够的CSMAR个股日行情，应保留原月频版本另建验证文件，在相同名单下检查调仓前后价格、分红复权、停复牌和持有收益；再比较月度等权与季度权重漂移。这样可以先判断收益实现差异，再判断选股差异，避免一次改变全部环节而无法解释原因。目前这些工作仍是待执行事项，正文没有将其写成已完成的证明。

最后才是在冻结规则下评价新增数据。全市场排名与候选池排名两条路径都应保留，风控分别与各自底层配对；任何核心优先、负PE排名移动或参数新设定都应另存文件、登记具体变化，不覆盖现有策略。这样才能继续保持本项目从基础到延展、再到优化的可读性和可审核性。
